\documentclass[10pt,a4paper]{amsart}
\usepackage{amsmath,amssymb,amsthm,graphicx}
% For commutative diagrams we use Xy-pic
\usepackage[all]{xy}

% We use 2cell for 2-commutative diagrams.
\xyoption{2cell}
% Theorem environments.
%
\theoremstyle{plain}
\newtheorem{theorem}[subsection]{Theorem}
\newtheorem{proposition}[subsection]{Proposition}
\newtheorem{lemma}[subsection]{Lemma}

\theoremstyle{definition}
\newtheorem{definition}[subsection]{Definition}
\newtheorem{example}[subsection]{Example}
\newtheorem{exercise}[subsection]{Exercise}
\newtheorem{situation}[subsection]{Situation}

\theoremstyle{remark}
\newtheorem{remark}[subsection]{Remark}
\newtheorem{remarks}[subsection]{Remarks}

\numberwithin{equation}{subsection}

% Macros
%
\def\lim{\mathop{\mathrm{lim}}\nolimits}
\def\colim{\mathop{\mathrm{colim}}\nolimits}
\def\Spec{\mathop{\mathrm{Spec}}}
\def\Hom{\mathop{\mathrm{Hom}}\nolimits}
\def\Ext{\mathop{\mathrm{Ext}}\nolimits}
\def\SheafHom{\mathop{\mathcal{H}\!\mathit{om}}\nolimits}
\def\SheafExt{\mathop{\mathcal{E}\!\mathit{xt}}\nolimits}
\def\Sch{\mathit{Sch}}
\def\Mor{\mathop{\mathrm{Mor}}\nolimits}
\def\Ob{\mathop{\mathrm{Ob}}\nolimits}
\def\Sh{\mathop{\mathit{Sh}}\nolimits}
\def\NL{\mathop{N\!L}\nolimits}
\def\CH{\mathop{\mathrm{CH}}\nolimits}
\def\proetale{{pro\text{-}\acute{e}tale}}
\def\etale{{\acute{e}tale}}
\def\QCoh{\mathit{QCoh}}
\def\Ker{\mathop{\mathrm{Ker}}}
\def\Im{\mathop{\mathrm{Im}}}
\def\Coker{\mathop{\mathrm{Coker}}}
\def\Coim{\mathop{\mathrm{Coim}}}

% Boxtimes
%
\DeclareMathSymbol{\boxtimes}{\mathbin}{AMSa}{"02}

%
% Macros for moduli stacks/spaces
%
\def\QCohstack{\mathcal{QC}\!\mathit{oh}}
\def\Cohstack{\mathcal{C}\!\mathit{oh}}
\def\Spacesstack{\mathcal{S}\!\mathit{paces}}
\def\Quotfunctor{\mathrm{Quot}}
\def\Hilbfunctor{\mathrm{Hilb}}
\def\Curvesstack{\mathcal{C}\!\mathit{urves}}
\def\Polarizedstack{\mathcal{P}\!\mathit{olarized}}
\def\Complexesstack{\mathcal{C}\!\mathit{omplexes}}
% \Pic is the operator that assigns to X its picard group, usage \Pic(X)
% \Picardstack_{X/B} denotes the Picard stack of X over B
% \Picardfunctor_{X/B} denotes the Picard functor of X over B
\def\Pic{\mathop{\mathrm{Pic}}\nolimits}
\def\Picardstack{\mathcal{P}\!\mathit{ic}}
\def\Picardfunctor{\mathrm{Pic}}
\def\Deformationcategory{\mathcal{D}\!\mathit{ef}}
\usepackage{cleveref}
\crefname{originaldefinition}{original definition}{original definitions}
\Crefname{originaldefinition}{Original definition}{Original definitions}
\makeatletter
\newlabel{situation-setup}{{02QL}{source}{}{sourcepointer-1}{}}
\newlabel{situation-setup@cref}{{[originaldefinition][0][]02QL}{[1][0][]source}{}{}{}}
\makeatother
\crefname{originallemma}{original lemma}{original lemmas}
\Crefname{originallemma}{Original lemma}{Original lemmas}
\makeatletter
\newlabel{lemma-chern-classes-tensor-product}{{0FA6}{source}{}{sourcepointer-2}{}}
\newlabel{lemma-chern-classes-tensor-product@cref}{{[originallemma][0][]0FA6}{[1][0][]source}{}{}{}}
\makeatother
\crefname{originallemma}{original lemma}{original lemmas}
\Crefname{originallemma}{Original lemma}{Original lemmas}
\makeatletter
\newlabel{lemma-chern-character-multiplicative}{{0F9D}{source}{}{sourcepointer-3}{}}
\newlabel{lemma-chern-character-multiplicative@cref}{{[originallemma][0][]0F9D}{[1][0][]source}{}{}{}}
\makeatother
\crefname{originalremark}{original remark}{original remarks}
\Crefname{originalremark}{Original remark}{Original remarks}
\makeatletter
\newlabel{remark-splitting-principle-perfect}{{0FAD}{source}{}{sourcepointer-4}{}}
\newlabel{remark-splitting-principle-perfect@cref}{{[originalremark][0][]0FAD}{[1][0][]source}{}{}{}}
\makeatother
\title{Chern classes, Chern character and Chern polynomials of the derived tensor product of two perfect objects whose Chern classes are defined}
\author{The Stacks Project Authors}
\date{stacks-project:chow.tex:local-snapshot-20261012; GNU FDL 1.2 or later}
\begin{document}
\maketitle
\noindent\textit{Editorial note.} All delivered mathematical text is the author's own, copied from the pinned Stacks Project chapter file chow.tex; no statement or proof step was rewritten, reordered or removed.\par\medskip
\noindent\textit{Editorial note.} Preamble provenance: the pinned chapter file opens with an input of the shared preamble and carries no preamble text of its own, so the project's single Stacks-wide preamble staged at sources/stacks-chow/preamble.tex is excerpted verbatim as two spans, lines 14-18 (Xy-pic setup) and lines 154-217 (the native theorem declarations and mathematical macros). Nothing was added, removed or redefined; the wrapper class and the amsmath/amssymb/amsthm/graphicx base packages are supplied by the assembler.\par\medskip
\noindent\textit{Editorial note.} The author closes the c\_2 computation with a sentence left to the reader; it is kept verbatim in the delivered proof (source line 9675): "which the reader can verify agrees with the formula in the statement of the lemma up to a factor of 2". This is the only step that is not written out; the binomial identity for P\_p and the c\_1 identity are fully proven inside the counted span.\par\medskip
\noindent\textit{Editorial note.} Canonical location of the counted proposition: Stacks tag 0FAG (sources/stacks-chow/pinned-original-source.tex line 9601).\par\medskip
\noindent\textit{Allowed standard prerequisites (editorial).}\begin{itemize}
\item Chern classes of a perfect object and the polynomials P\_p with r(E) = P\_0(E). Same chapter, chow.tex, Remark labelled remark-splitting-principle-perfect (source line 9494) and the surrounding definition of Chern classes of perfect objects. Supplies the relation between the P\_p and the Chern classes c\_i used at proof lines 9657-9673, and the splitting principle invoked at proof line 9655.
\item The corresponding tensor product and Chern character statements for vector bundles. Same chapter, Lemma labelled lemma-chern-classes-tensor-product (source line 8567) and Lemma labelled lemma-chern-character-multiplicative (source line 8986), invoked at proof lines 9628-9630. Supply the vector bundle case which the envelope reduction would reduce the lemma to; the combinatorial P\_p verification in the counted span is the author's preferred second route.
\item The base situation (S, delta) for bivariant operational Chow groups. Same chapter, Situation labelled situation-setup (source line 1372), invoked in the statement. Fixes the ambient situation and the ring A\textasciicircum{}*(X) in which the identities are stated; no further step of the proof uses it.
\end{itemize}\medskip
\ifdefined\hypertarget\hypertarget{sourcepointer-1}{}\fi
\noindent\textit{Source reference.} Original definition 02QL: sources/stacks-chow/pinned-original-source.tex, source lines 1372--1372. Same-chapter pointer into the author environment situation at source line 1372; invoked at proof line None of the counted proof. It is retained only as a navigation pointer; every load-bearing step is inside the delivered span.\par\medskip
\ifdefined\hypertarget\hypertarget{sourcepointer-2}{}\fi
\noindent\textit{Source reference.} Original lemma 0FA6: sources/stacks-chow/pinned-original-source.tex, source lines 8567--8567. Same-chapter pointer into the author environment lemma at source line 8567; invoked at proof line 9629 of the counted proof. It is retained only as a navigation pointer; every load-bearing step is inside the delivered span.\par\medskip
\ifdefined\hypertarget\hypertarget{sourcepointer-3}{}\fi
\noindent\textit{Source reference.} Original lemma 0F9D: sources/stacks-chow/pinned-original-source.tex, source lines 8986--8986. Same-chapter pointer into the author environment lemma at source line 8986; invoked at proof line 9630 of the counted proof. It is retained only as a navigation pointer; every load-bearing step is inside the delivered span.\par\medskip
\ifdefined\hypertarget\hypertarget{sourcepointer-4}{}\fi
\noindent\textit{Source reference.} Original remark 0FAD: sources/stacks-chow/pinned-original-source.tex, source lines 9494--9494. Same-chapter pointer into the author environment remark at source line 9494; invoked at proof line 9632 of the counted proof. It is retained only as a navigation pointer; every load-bearing step is inside the delivered span.\par\medskip
\begin{lemma}
\label{lemma-chern-classes-perfect-tensor-product}
In Situation \ref{situation-setup} let $X$ be locally of finite type over $S$.
Let $E$ and $F$ be perfect objects of $D(\mathcal{O}_X)$ whose Chern classes
are defined. Then we have
$$
c_1(E \otimes_{\mathcal{O}_X}^\mathbf{L} F) =
r(E) c_1(\mathcal{F}) + r(F) c_1(\mathcal{E})
$$
and for $c_2(E \otimes_{\mathcal{O}_X}^\mathbf{L} F)$ we have the expression
$$
r(E) c_2(F) + r(F) c_2(E) + {r(E) \choose 2} c_1(F)^2 +
(r(E)r(F) - 1) c_1(F)c_1(E) + {r(F) \choose 2} c_1(E)^2
$$
and so on for higher Chern classes in $A^*(X)$. Similarly, we have
$ch(E \otimes_{\mathcal{O}_X}^\mathbf{L} F) = ch(E) ch(F)$
in $A^*(X) \otimes \mathbf{Q}$. More precisely, we have
$$
P_p(E \otimes_{\mathcal{O}_X}^\mathbf{L} F) = \sum\nolimits_{p_1 + p_2 = p}
{p \choose p_1} P_{p_1}(E) P_{p_2}(F)
$$
in $A^p(X)$.
\end{lemma}

\begin{proof}
After choosing an envelope $f : Y \to X$ such that $Lf^*E$ and $Lf^*F$
can be represented by locally bounded complexes of finite locally
free $\mathcal{O}_X$-modules this follows by a computation from the
corresponding result for vector bundles in
Lemmas \ref{lemma-chern-classes-tensor-product} and
\ref{lemma-chern-character-multiplicative}.
A better proof is probably to use the splitting principle as in
Remark \ref{remark-splitting-principle-perfect}
and reduce the lemma to computations in polynomial rings 
which we describe in the next paragraph.

\medskip\noindent
Let $A$ be a commutative ring (for us this will be the subring of the
bivariant chow ring of $X$ generated by Chern classes).
Let $S$ be a finite set together with maps $\epsilon : S \to \{\pm 1\}$
and $f : S \to A$. Define
$$
P_p(S, f , \epsilon) = \sum\nolimits_{s \in S} \epsilon(s) f(s)^p
$$
in $A$. Given a second triple $(S', \epsilon', f')$
the equality that has to be shown for $P_p$ is the equality
$$
P_p(S \times S', f + f' , \epsilon \epsilon') = 
\sum\nolimits_{p_1 + p_2 = p}
{p \choose p_1} P_{p_1}(S, f, \epsilon) P_{p_2}(S', f', \epsilon')
$$
To see this is true, one reduces to the polynomial ring on variables
$S \amalg S'$ and one shows that each term $f(s)^if'(s')^j$ occurs
on the left and right hand side with the same coefficient.
To verify the formulas for $c_1(E \otimes_{\mathcal{O}_X}^\mathbf{L} F)$
and $c_2(E \otimes_{\mathcal{O}_X}^\mathbf{L} F)$ we use the splitting
principle to reduce to checking these formulae in a torsion free ring.
Then we use the relationship between $P_j(E)$ and $c_i(E)$ proved
in Remark \ref{remark-splitting-principle-perfect}. For example
$$
c_1(E \otimes F) = P_1(E \otimes F) = r(F)P_1(E) + r(E)P_1(F) =
r(F)c_1(E) + r(E)c_1(F)
$$
the middle equation because $r(E) = P_0(E)$ by definition. Similarly, we have
\begin{align*}
& 2c_2(E \otimes F) \\
& = c_1(E \otimes F)^2 - P_2(E \otimes F) \\
& =
(r(F)c_1(E) + r(E)c_1(F))^2 -
r(F)P_2(E) - P_1(E)P_1(F) - r(E)P_2(F) \\
& =
(r(F)c_1(E) + r(E)c_1(F))^2 -
r(F)(c_1(E)^2 - 2c_2(E)) - c_1(E)c_1(F) - \\
& \quad r(E)(c_1(F)^2 - 2c_2(F))
\end{align*}
which the reader can verify agrees with the formula in the statement
of the lemma up to a factor of $2$.
\end{proof}

\end{document}
